\documentclass[../../../include/open-logic-section]{subfiles}

\begin{document}

\iftag{FOL}
      {\olfileid{fol}{ntd}{ppr}}
      {\olfileid{pl}{ntd}{ppr}}

\olsection{可推导性与命题联结词}

\begin{explain}
我们将证明，自然演绎的可推导关系~$\Proves$ 足够强，
能够确立一些涉及命题联结词的基本事实，
例如 $!A \land !B \Proves !A$ 和 $!A, !A \lif !B \Proves !B$（肯定前件式）。
在完备性定理的证明中需要用到这些事实。
\end{explain}

% Part: first-order-logic
% Chapter: natural-deduction
% Section: provability-propositional

% verification of properties of provability needed for maximally
% consistent sets in the completeness chapter.

\begin{prop}\ollabel{prop:provability-land}
  \begin{enumerate}
  \item \ollabel{prop:provability-land-left} $!A \land !B \Proves !A$ 且 $!A \land !B \Proves !B$。
  \item \ollabel{prop:provability-land-right} $!A, !B \Proves !A \land !B$。
  \end{enumerate}
\end{prop}

\begin{proof}
  \begin{enumerate}
  \item 我们可以分别推导如下：
    \begin{prooftree}
      \AxiomC{$!A \land !B$}
      \RightLabel{\Elim{\land}}
      \UnaryInfC{$!A$}
      \DisplayProof\qquad\bottomAlignProof
      \AxiomC{$!A \land !B$}
      \RightLabel{\Elim{\land}}
      \UnaryInfC{$!B$}
    \end{prooftree}
  \item 我们可以推导如下：
    \begin{prooftree}
      \AxiomC{$!A$}
      \AxiomC{$!B$}
      \RightLabel{\Intro{\land}}
      \BinaryInfC{$!A \land !B$}
    \end{prooftree}
  \end{enumerate}
\end{proof}

\begin{prop}\ollabel{prop:provability-lor}
  \begin{enumerate}
  \item $!A \lor !B, \lnot !A, \lnot !B$ 是不一致的。
  \item $!A \Proves !A \lor !B$ 且 $!B \Proves !A \lor !B$。
  \end{enumerate}
\end{prop}

\begin{proof}
  \begin{enumerate}
  \item 考虑如下推导：
    \begin{prooftree}
      \AxiomC{$!A \lor !B$}
      \AxiomC{$\lnot !A$}
      \AxiomC{$\Discharge{!A}{1}$}
      \RightLabel{\Elim{\lnot}}
      \BinaryInfC{$\lfalse$}
      \AxiomC{$\lnot !B$}
      \AxiomC{$\Discharge{!B}{1}$}
      \RightLabel{\Elim{\lnot}}
      \BinaryInfC{$\lfalse$}
      \DischargeRule{\Elim{\lor}}{1}
      \TrinaryInfC{$\lfalse$}
    \end{prooftree}
    这是一个从未解除的假设 $!A \lor !B$、$\lnot !A$ 和 $\lnot !B$ 推出 $\lfalse$ 的推导。
  \item 我们可以分别推导如下：
    \begin{prooftree}
      \AxiomC{$!A$}
      \RightLabel{\Intro{\lor}}
      \UnaryInfC{$!A \lor !B$}
      \DisplayProof\qquad\bottomAlignProof
      \AxiomC{$!B$}
      \RightLabel{\Intro{\lor}}
      \UnaryInfC{$!A \lor !B$}
    \end{prooftree}
  \end{enumerate}
\end{proof}

\begin{prop}\ollabel{prop:provability-lif}
  \begin{enumerate}
  \item \ollabel{prop:provability-lif-left}  $!A, !A \lif !B \Proves !B$。
  \item \ollabel{prop:provability-lif-right}
    $\lnot !A \Proves !A \lif !B$ 且 $!B \Proves !A \lif !B$。
  \end{enumerate}
\end{prop}

\begin{proof}
  \begin{enumerate}
  \item 我们可以推导如下：
    \begin{prooftree}
      \AxiomC{$!A \lif !B$}
      \AxiomC{$!A$}
      \RightLabel{\Elim{\lif}}
      \BinaryInfC{$!B$}
    \end{prooftree}

  \item 这可通过以下两个推导展示：
    \begin{prooftree}
      \AxiomC{$\lnot !A$}
      \AxiomC{$\Discharge{!A}{1}$}
      \RightLabel{\Elim{\lnot}}
      \BinaryInfC{$\lfalse$}
      \RightLabel{\FalseInt}
      \UnaryInfC{$!B$}
      \DischargeRule{\Intro{\lif}}{1}
      \UnaryInfC{$!A \lif !B$}
      \DisplayProof\qquad\bottomAlignProof
      \AxiomC{$!B$}
      \RightLabel{\Intro{\lif}}
      \UnaryInfC{$!A \lif !B$}
    \end{prooftree}
    注意，$\Intro{\lif}$ 可以解除假设~$!A$，但并非必须解除。
  \end{enumerate}
\end{proof}

\end{document}